% A.79: full root, DE and PBC proof/source audits PASS.
\subsection{Cyclic-derived closures at a supersingular place}

\begin{statement}{A.79}{Proved}{Cyclic-derived prime-power closures and coprime abelian factors}
Let \(p=2\) or let \(p\) be an odd prime with
\(p\equiv3\pmod4\). Let \(P\) be a finite
\(p\)-group with cyclic derived subgroup, and let
\(B\) be a finite abelian group of order prime to \(p\).
Every actual connected finite \'etale cover of
\(E\setminus\{O\}\) whose geometric Galois closure
is a quotient of \(P\times B\) satisfies both assertions
of Task A.5 at every coefficient prime. Every actual
intermediate, original induced cusp packet and permitted
open derived subgroup is included.

The reduction place in this argument is a supersingular
place of the fixed CM elliptic curve. The conclusion
neither treats arbitrary nilpotent closures nor asserts
cofinality. In particular, the cyclic-derived hypothesis
supplies an actual cusp-divisor relation which the rank-zero
theorem A.78 alone does not supply for an arbitrary group.
\end{statement}

\begin{proof}
\paragraph{The actual prime-power closure and its inertia.}
First suppose the actual Galois closure has group \(P_0\),
a quotient of \(P\). Its derived subgroup remains cyclic.
The two actual handle generators \(a,b\) generate \(P_0\).
Their commutator \(c=[a,b]\) normally generates \(P_0'\),
because the quotient killing it is abelian. Since every
subgroup of cyclic \(P_0'\) is characteristic in
\(P_0'\), the subgroup \(\langle c\rangle\) is
normal in \(P_0\). Hence
\[
                   I=\langle c\rangle=P_0'.             \tag{A.79.1}
\]
This is the actual geometric cusp inertia. In particular
it is normal. No inertia is inferred from an abstract
curve isomorphism or a dimension count.

Let \(C\to E\) be the proper closure. The actual
quotient \(E_1=C/I\to E\) is unramified on the
proper curves, and becomes an elliptic isogeny after a
choice of origin. Its kernel is \(P_0/I\). Write
\(d\) for the exponent of this kernel. Each point of
the kernel has a unique cusp of \(C\) above it,
with ramification index \(|I|\). If \(Q,Q'\) are
kernel points and \(P_Q,P_{Q'}\) their cusps, then
pullback of the actual principal elliptic divisor gives
\[
 d(Q-Q')\sim0
 \quad\Longrightarrow\quad
 d|I|(P_Q-P_{Q'})\sim0.                                \tag{A.79.2}
\]
Thus every cusp difference is killed by a power of \(p\).

At the chosen place above \(p\), the fixed elliptic
curve has potential good supersingular reduction. After
a finite local extension A.78 therefore applies to the
actual \(p\)-group cover, whose only branch section
is \(O\). It gives potential good reduction of \(J_C\)
of special \(p\)-rank zero.

\paragraph{Every actual intermediate retains companions.}
For an actual subgroup \(H<P_0\), let \(Z=C/H\).
The norm \(J_C\to J_Z\) is surjective, so potential
good reduction and special rank zero descend to \(J_Z\).
Over an algebraic closure the pushforward of a point
under \(C\to Z\) has residue degree one. Pushing
forward (A.79.2) therefore proves the same torsion
relation for any two cusps of \(Z\); no ramification
factor is inserted into this pushforward.

Since \(I\) is normal, the number of cusps of \(Z\)
is \([P_0:HI]\). Every proper subgroup of a finite
\(p\)-group is contained in a maximal subgroup, and
\(P_0'\subset\Phi(P_0)\) lies in each such maximal
subgroup. Thus \(HI\ne P_0\), and the number of
cusps is at least \(p\). The compact genus is at
least one, because there is a finite nonconstant map to
\(E\). If it is one, A.40 already applies.
Otherwise the rank-zero abelian special fibre has no
nonzero geometric \(p\)-power torsion. All cusp differences
therefore specialize to zero in the good Jacobian.
The compact-type Abel-map and rational-path argument of
A.41 identifies their marked joining paths as entirely
rational, including the marked endpoint components.
Every cusp has a distinct companion. A.42 gives all-prime
fixed inertia, and A.59 gives all normalizers. The base
quotient \(Z=E\), corresponding to \(H=P_0\), and every abelian \(P_0\)
case are the genus-one assertion A.40.

\paragraph{An actual connected coprime abelian pullback.}
Now allow the factor \(B\). Every subgroup of a direct
product of finite groups of coprime orders splits as the
product of its intersections with the factors: Bezout
powers extract \((x,1)\) and \((1,y)\) from \((x,y)\).
This applies to the normal quotient kernel and the actual
covering subgroup. Consequently the actual closure group
is \(P_0\times B_0\), with \(P_0'\) cyclic and
\(B_0\) abelian of order prime to \(p\), and its
covering subgroup is \(H_p\times H_B\).

The proper \(B_0/H_B\)-cover is an actual elliptic
isogeny \(\eta:E'\to E\) of degree prime to \(p\):
the peripheral commutator vanishes in the abelian quotient.
Let \(Z\to E\) be the actual \(H_p\)-intermediate
of the \(P_0\)-closure. Its degree is a power of \(p\).
The isogeny extension of function fields is Galois of
coprime degree, so these two fields are linearly disjoint.
Their actual compositum is therefore the field of the
geometrically connected fibre product
\[
                              Y=Z\times_E E'.           \tag{A.79.3}
\]
Proper smooth normalization is understood. The removed
points are exactly the pairs of a cusp of \(Z\) and
a point of \(\ker\eta\). If \(Z\) has genus one,
this is an unramified proper elliptic cover and A.40
applies. Otherwise the preceding argument supplies the
actual rational companion paths on \(Z\).

\paragraph{Extension of the map and a separate compact-type argument.}
Choose a finite extension defining all algebraic maps,
marked points and good reduction data. On the compact-type
marked model \(\mathcal Z\), extend the Abel map based
at a cusp using \cite[Theorem 1.1 and Proposition 3.2]{fghDupuyRabinoff}.
Compose with the actual homomorphism \(J_Z\to E\),
which extends between the good abelian schemes. This gives
an actual morphism \(\mathcal Z\to\mathcal E\)
across the nodes and the inserted rational components.
The prime-to-\(p\) isogeny extends to a finite \'etale
isogeny of the good elliptic models. Thus
\[
 \mathcal Y=\mathcal Z\times_{\mathcal E}\mathcal E'
                                                               \tag{A.79.4}
\]
is a proper flat semistable model of the actual connected
curve, finite \'etale over \(\mathcal Z\).

Its graph is not declared a tree merely from finite
\'{e}taleness. Instead apply Lehr--Matignon Theorem 6.13
\cite[Theorem 6.13]{fghLehrMatignon2005} to the actual nilpotent Galois closure over the good
elliptic model, with its sole \'etale branch section
\(O\). The closure Jacobian has potential good reduction,
and its norm quotient \(J_Y\) does too. Hence the
semistable dual graph of \(\mathcal Y\) is a tree.
No special-rank-zero assertion for \(J_Y\) is needed.

\paragraph{The selected companion path lifts as a literal copy.}
Let \(T\) be the rational subtree joining a cusp
\(P\) of \(Z\) to a distinct companion \(Q\).
Every component of \(T\) maps constantly to the good
elliptic special fibre, and those constants agree at the
nodes. Its endpoint sections map to \(O\), so its
whole image is \(O\). The fibre of the good \'etale
isogeny above \(O\) is a reduced discrete set of
\(\deg\eta\) points. Its pullback over \(T\) is
literally that many disjoint copies of \(T\).

A selected cusp of \(Y\) chooses one copy. In that
copy the lift of \(Q\) is a distinct marked point,
with the same kernel coordinate. Distinct prime-to-\(p\)
kernel sections have distinct reductions. The joining
path remains rational after marked stabilization, and
the separately verified tree property makes it the
unique path. This verifies the exact hypothesis of A.42
for each selected cusp of \(Y\). A.42 proves all
fixed-inertia conclusions and A.59 supplies normalizers
for the nilpotent closure.

This is the actual preservation argument of A.66/A.76
with residue prime \(p\) in place of two. Its changed
hypotheses have all been verified above: good reduction
of the elliptic target, prime-to-\(p\) good isogeny,
connected actual compositum, compact type independently,
and literal companion-path copies. Finite extensions
are used only to define the algebraic and reduction data.
The arithmetic comparisons are up to inner conjugacy,
with the established restriction to the permitted open
derived subgroups; no geometric path is asserted to be
defined over a finite field extension. Thus all original
cusp packets are retained. No quotient fixed element is
lifted, and no cofinality conclusion is asserted.
\end{proof}
